%! AP Statistics — Unit 2, Lesson 2.5 — Exit Ticket (Answer Key)
\documentclass[10pt]{article}
\usepackage{apstats-article}
\usepackage{apstats-key}

\begin{document}

\pageheader{Unit 2, Lesson 2.5}{Exit Ticket --- Answer Key}
\namedateperiod

\vspace{0.18in}

\begin{tcolorbox}[colback=sky, colframe=navy, boxrule=0.9pt, arc=2mm,
    left=4mm, right=4mm, top=3mm, bottom=3mm]
\small\textit{A consumer testing agency measured engine size (in liters) and highway fuel
efficiency (in mpg) for a random sample of 30 cars. The correlation coefficient was
$r = -0.87$.}
\end{tcolorbox}

\vspace{0.2in}

\begin{enumerate}[leftmargin=1.6em, itemsep=0.22in]

  \item Write a \textbf{complete interpretation} of $r = -0.87$ in context. Your response
        must include the direction, strength, and a statement that the association is linear.\\[10pt]
        \ansline{There is a strong negative linear association between engine size (liters) and
        highway fuel efficiency (mpg) for these 30 cars ($r = -0.87$). As engine size increases,
        fuel efficiency tends to decrease.}

  \item A classmate says, ``Since $r$ is close to $-1$, we can conclude that having a
        larger engine \emph{causes} a car to get lower gas mileage.'' Is your classmate
        correct? Explain.\\[10pt]
        \ansline{No. Correlation does not imply causation. A strong $r$ shows the two variables
        are associated, but does not prove that one causes the other. Other variables --- such as
        vehicle weight or aerodynamics --- may explain the relationship.}

\end{enumerate}

\vspace{0.14in}

\begin{tcolorbox}[colback=redbg, colframe=redacc, boxrule=0.8pt, arc=2mm,
    left=4mm, right=4mm, top=3mm, bottom=3mm]
\small\textbf{AP-Style Practice}\\[6pt]
A researcher finds that the correlation between two variables is $r = 0.95$. Which of the
following conclusions is best supported?

\vspace{10pt}
\begin{enumerate}[label=\textbf{(\Alph*)}, leftmargin=2em, itemsep=6pt]
  \item \ans{\textbf{(A)}} There is a strong positive \emph{linear} association between the two variables.
  \item An increase in one variable causes an increase in the other variable.
  \item The relationship between the two variables is definitely linear.
  \item There is no need to look at the scatterplot because $r$ is near 1.
\end{enumerate}
\end{tcolorbox}

\vspace{0.1in}
\noindent\textcolor{slate}{\small My answer: \ans{A} \quad Because: \ans{$r$ measures the
strength and direction of a \emph{linear} association; it does not prove causation, guarantee
linearity, or make the scatterplot unnecessary.}}

\begin{teachernote}
\textbf{(B)} is wrong because correlation never proves causation. \textbf{(C)} is wrong because
a near-perfect $r$ could occur with a consistently rising curve. \textbf{(D)} is wrong because
you should always examine the scatterplot --- Anscombe's Quartet famously shows four very
different datasets with the same $r$.
\end{teachernote}

\end{document}
